% TEMPLATE-MILC-v3.tex - plantilla maestra UNICA de las guias MILC v3.
%
% La usan las cuatro familias de guias, cada una con su driver:
%   · Grados 6-11          -> scripts/build-guias-g11.py
%   · Bebras               -> scripts/build-guias-bebras.py
%   · Semillero            -> scripts/build-guias-semillero.py
%   · Territorio interior  -> scripts/build-guias-territorio-interior.py
%
% Antes habia tres copias casi identicas (~400 lineas iguales, 6 distintas):
% cada mejora habia que aplicarla tres veces, y en la practica no se hacia
% (el motor de recursos vivia solo en dos de ellas). Lo que varia entre
% familias esta parametrizado en 6 placeholders que cada driver llena
% (se nombran aqui SIN su sintaxis real: el builder reemplaza por texto plano,
% tambien dentro de los comentarios, y un valor multilinea rompe el preambulo):
%
%   HEADER_IZQ        encabezado izquierdo de pagina
%   HEADER_DER        encabezado derecho de pagina
%   PORTADA_ETIQUETA  rotulo sobre el numero grande (GRADO/MOMENTO/MODULO)
%   PORTADA_SERIE     linea de serie de la portada
%   PORTADA_ID        identificador de la guia en la portada
%   PORTADA_CIRCULO   el circulito de grado (°), o vacio si no aplica
%   PDF_TITULO        titulo en texto plano (sin \textbf ni comillas TeX) para
%                     los metadatos del PDF (/Title); lo produce lib_xelatex.texto_pdf
%
% Composicion (auditoria 2026-09-03): las bandas de estacion piden espacio
% (\Needspace*) y prohiben el salto justo despues (\nopagebreak), asi no quedan
% huerfanas al pie; las cajas largas (softbox, drawbox, infoband, puentebox) son
% `breakable`, asi una caja de media pagina ya no arrastra todo a la siguiente
% dejando la anterior al 40 %. Todo texto sobre mostaza o verde va en milcNegro
% (blanco sobre #E5B400 da 1,93:1 y sobre #58A923 2,95:1; ninguno pasa WCAG AA).
% El resto de claves las documenta content/guias/_SCHEMA.md. El script valida
% que no queden placeholders sin reemplazar antes de escribir el archivo final.

% Etiquetado PDF (latex-lab): /Lang, estructura y texto alternativo de las
% figuras, para lectores de pantalla. Debe ir ANTES de \documentclass.
\DocumentMetadata{lang=es-CO,tagging=on}
\documentclass[11pt,letterpaper]{article}
% headheight=24pt: la cabecera derecha lleva dos lineas (identidad de la guia
% y estacion actual). headsep se acorta en la misma medida para que el bloque
% de texto quede exactamente donde estaba con headheight=12pt.
\usepackage[margin=2.0cm,headheight=24pt,headsep=13pt]{geometry}
\usepackage[table]{xcolor}
\usepackage{fontspec}
\usepackage[spanish]{babel}
\usepackage{tikz}
% Librerías TikZ para los diagramas de las guías (bloque `recursos:`):
%  · babel  → babel en español vuelve activos caracteres como «>», y TikZ los usa
%             en sus opciones (>=stealth, ->). Sin esta librería, cualquier
%             diagrama con flechas revienta con «\language@active@arg> has an
%             extra }». Debe ir ANTES de que se dibuje nada.
%  · positioning        → sintaxis «below left=2pt and 3pt».
%  · decorations.pathreplacing → llaves ({brace}) para señalar rangos.
\usetikzlibrary{babel,positioning,decorations.pathreplacing}
\usepackage{graphicx}
% Circuitos: el trabajo de electrónica se hace hoy con micro:bit y sus sensores
% integrados (MakeCode, sin cableado), así que no se necesitan esquemáticos.
% Si algún día entran montajes con protoboard, basta descomentar esta línea:
% los diagramas .tex/.tikz declarados en `recursos:` ya se incrustan solos.
% \usepackage{circuitikz}
\usepackage[most]{tcolorbox}
\usepackage{tabularx}
\usepackage{array}
\usepackage{fancyhdr}
\usepackage{titlesec}
\usepackage[document]{ragged2e}  % bandera a la derecha en todo el cuerpo
\usepackage{enumitem}
\usepackage{microtype}
\usepackage{etoolbox}
\usepackage{needspace}
% hyperref al final del preambulo: metadatos del PDF (titulo, autor, idioma,
% familia), marcadores por estacion (\pdfbookmark) y enlaces sin recuadro.
\usepackage[hidelinks,
  pdftitle={La rueda y el agua — máquinas simples del campo colombiano},
  pdfauthor={PhD. Álvaro Cárdenas Orozco},
  pdfsubject={Tecnología e Informática · Grado 9 · Guía 2 · MILC v3},
  pdflang={es-CO},
  bookmarksopen=true,bookmarksnumbered=false]{hyperref}

% Helvetica Neue no trae la flecha U+2192, asi que cada «→» escrito literalmente
% en el YAML se compilaba a NADA, en silencio (462 flechas en 61 guias: rutas del
% tipo «paleta → tipografia → lockup» salian rotas). Mapeamos el caracter a su
% version matematica, que si tiene glifo. Asi el autor puede seguir escribiendo
% «→» comodamente en el YAML.
\usepackage{newunicodechar}
\newunicodechar{→}{\ensuremath{\rightarrow}}
% Mismo problema con los simbolos de marca y vinetas: el «✓» de «una palomita ✓»
% salia como cuadrito vacio en el PDF de todas las guias que lo usaban.
\usepackage{pifont}
\newunicodechar{✓}{\ding{51}}
\newunicodechar{✗}{\ding{55}}
\newunicodechar{✔}{\ding{52}}
\newunicodechar{•}{\textbullet}
\newunicodechar{≥}{\ensuremath{\geq}}
\newunicodechar{≤}{\ensuremath{\leq}}

\setmainfont{Helvetica Neue}
\setsansfont{Helvetica Neue}
\setlength{\parindent}{0pt}
\setlength{\parskip}{6pt}
% Sin particion de palabras: en 6.º-7.º una palabra cortada por guion al final
% de linea («Comuni-cacion») frena la lectura mucho mas que una linea corta.
\hyphenpenalty=10000
\exhyphenpenalty=10000
\pretolerance=2000
\tolerance=3000
\setlength{\emergencystretch}{3em}
\linespread{1.10}
\renewcommand{\arraystretch}{1.25}
\setlist{leftmargin=5mm,itemsep=1mm,topsep=1mm}
\newcolumntype{Y}{>{\RaggedRight\arraybackslash}X}

\definecolor{milcMagenta}{HTML}{E6007E}
\definecolor{milcVino}{HTML}{5A0038}
\definecolor{milcTurquesa}{HTML}{0093A5}
\definecolor{milcVerde}{HTML}{58A923}
\definecolor{milcMostaza}{HTML}{E5B400}
\definecolor{milcOcre}{HTML}{B86600}
\definecolor{milcNegro}{HTML}{241816}
\definecolor{milcGris}{HTML}{F7F7F4}
\definecolor{milcLinea}{HTML}{E8E2DD}
\definecolor{milcGradePrimary}{HTML}{7C3AED}
\definecolor{milcGradeDark}{HTML}{4C1D95}
\definecolor{milcGradeSoft}{HTML}{EDE9FE}
\definecolor{milcGradeAccent}{HTML}{CFE7FF}
\definecolor{milcGradeText}{HTML}{FFFFFF}
\color{milcNegro}

\pagestyle{fancy}
\fancyhf{}
% Cabecera de dos lineas: identidad de la guia arriba y, a la derecha y en
% gris, la estacion en curso (\markright desde \milcestacion). Antes de la
% primera estacion la segunda linea va vacia; el \strut mantiene la altura
% para que la linea de arriba no baile entre paginas.
\lhead{\small Tecnología e Informática · Grado 9\\[-1pt]{\footnotesize\strut}}
\rhead{\small Guía 2 · MILC v3\\[-1pt]{\footnotesize\strut\color{milcNegro!65}\rightmark}}
\cfoot{\small\thepage}
\renewcommand{\headrulewidth}{0pt}

% Sobre mostaza (#E5B400) y verde (#58A923) el texto va en milcNegro; sobre el
% resto de la paleta, en blanco. Se usa dentro de fonttitle/fontupper, donde un
% \color posterior manda sobre coltitle.
\newcommand{\milctintasobre}[1]{%
  \ifboolexpr{test{\ifstrequal{#1}{milcMostaza}} or test{\ifstrequal{#1}{milcVerde}}}%
    {\color{milcNegro}}{\color{white}}}

\newtcolorbox{titlebox}[1]{
  enhanced,colback=#1,colframe=#1,arc=7mm,boxrule=0pt,
  left=7mm,right=7mm,top=3mm,bottom=3mm,
  before skip=8pt,after skip=8pt,
  fontupper=\sffamily\bfseries\Large\color{white}
}

\newtcolorbox{softbox}[3]{
  enhanced,breakable,pad at break=3mm,
  colback=#2,colframe=#1,borderline west={4pt}{0pt}{#1},
  arc=2mm,boxrule=.4pt,left=4mm,right=4mm,top=3mm,bottom=3mm,
  before skip=6pt,after skip=8pt,title={#3},coltitle=white,
  colbacktitle=#1,fonttitle=\sffamily\bfseries\milctintasobre{#1},fontupper=\RaggedRight,
  attach boxed title to top left={xshift=3mm,yshift=-2mm},
  boxed title style={arc=2mm, boxrule=0pt}
}

\newtcolorbox{drawbox}[2]{
  enhanced,breakable,pad at break=4mm,
  colback=white,colframe=#1,arc=4mm,boxrule=1pt,
  left=5mm,right=5mm,top=4mm,bottom=4mm,before skip=8pt,after skip=8pt,
  title={#2},coltitle=white,colbacktitle=#1,fonttitle=\sffamily\bfseries\milctintasobre{#1},
  fontupper=\RaggedRight,
  attach boxed title to top left={xshift=4mm,yshift=-2mm},
  boxed title style={arc=3mm, boxrule=0pt}
}

\newtcolorbox{infoband}[1]{
  enhanced,breakable,pad at break=2mm,colback=#1!8,colframe=#1!50,arc=2mm,boxrule=.5pt,
  left=4mm,right=4mm,top=2mm,bottom=2mm,before skip=4pt,after skip=4pt,
  fontupper=\small\RaggedRight
}

% Puente narrativo entre secciones (contrato editorial Conéctate).
% Da continuidad: "Acabas de X. Ahora vas a Y porque Z."
% Visualmente: banda izquierda en color del grado, texto en itálica pequeña.
\newtcolorbox{puentebox}{
  enhanced,breakable,pad at break=2mm,colback=milcGradeSoft,colframe=milcGradeDark,
  borderline west={3pt}{0pt}{milcGradeDark},
  arc=0mm,boxrule=0pt,left=5mm,right=4mm,top=2.5mm,bottom=2.5mm,
  before skip=4pt,after skip=8pt,
  fontupper=\itshape\small\color{milcNegro}\RaggedRight
}
% La flecha va en modo matematico a proposito: Helvetica Neue NO tiene el glifo
% U+2192, asi que \textrightarrow se compilaba a NADA («Missing character: There
% is no → in font Helvetica Neue») y el puente salia sin flecha. En modo
% matematico la flecha viene de la fuente matematica y si se dibuja.
\newcommand{\puente}[1]{%
  \begin{puentebox}%
    \textcolor{milcGradeDark}{$\rightarrow$}\hspace{2mm}#1%
  \end{puentebox}%
}

\newcommand{\linead}[1]{\noindent\color{milcLinea}\rule{#1}{0.5pt}\color{milcNegro}}
\newcommand{\respuesta}[1]{\par\vspace{2mm}\foreach \n in {1,...,#1}{\noindent\color{milcLinea}\rule{\linewidth}{0.5pt}\par\vspace{5mm}}\color{milcNegro}}
\newcommand{\checkbox}{\textcolor{milcGradeDark}{\fbox{\phantom{x}}}\hspace{5pt}}

% ─── Listas aireadas para las guías socioemocionales ────────────────────
% Componer las verificaciones como «\checkbox texto\\» deja el texto que salta
% de línea debajo del cuadrito y apelmaza el bloque. Estos entornos alinean la
% continuación y airean los ítems. Los usa Territorio Interior; cualquier
% familia puede hacerlo.
\newlist{checklista}{itemize}{1}
\setlist[checklista]{
  label=\textcolor{milcGradeDark}{\fbox{\phantom{x}}},
  leftmargin=9mm, labelsep=3.2mm, itemsep=2.4mm,
  topsep=2.5mm, parsep=0pt, partopsep=0pt, before=\RaggedRight
}
\newlist{pasoslista}{enumerate}{1}
\setlist[pasoslista]{
  label=\textcolor{milcGradeDark}{\sffamily\bfseries\arabic*.},
  leftmargin=8mm, labelsep=3mm, itemsep=2.8mm,
  topsep=1mm, parsep=0pt, partopsep=0pt, before=\RaggedRight
}

\newcommand{\phasecard}[4]{
\begin{tcolorbox}[enhanced,colback=#3,colframe=#2,arc=3mm,boxrule=.5pt,height=3.05cm,valign=center,halign=center,left=1.5mm,right=1.5mm,top=2mm,bottom=2mm]
{\sffamily\bfseries\fontsize{22}{24}\selectfont\textcolor{#2}{#1}}\par
{\sffamily\bfseries\small #4}
\end{tcolorbox}
}

% ── Piezas del rediseno 2026 (legibilidad 6.º-11.º) ───────────────────
% Banda de estacion: reemplaza el titlebox macizo. Titulo a la izquierda,
% dato practico (tiempo/modalidad) a la derecha, en una sola linea.
% Nunca huerfana: pide 9 lineas libres (o salta de pagina rellenando con
% \vfill) y prohibe el salto justo despues de la banda. Nueve y no seis:
% la caja partible que sigue necesita ~80pt (titulo adosado, relleno y dos
% lineas) para arrancar en la misma pagina; con menos, tcolorbox la manda
% entera a la siguiente y la banda se queda sola. El penalty 10000 va
% en `after`, ANTES del espacio posterior: un `after skip` (aunque sea 0pt)
% es un \vskip tras la caja, y ese glue es un punto de corte legal que un
% \nopagebreak escrito despues de \end{tcolorbox} ya no cierra. Cada banda
% deja marcador PDF y actualiza la cabecera derecha.
\newcounter{milcestacion}
\newcommand{\milcestacion}[2]{%
\Needspace*{9\baselineskip}%
\stepcounter{milcestacion}%
\pdfbookmark[1]{#2}{milcestacion\themilcestacion}%
\markright{#2}%
\begin{tcolorbox}[enhanced,colback=#1,colframe=#1,arc=2mm,boxrule=0pt,
  left=5mm,right=5mm,top=2.6mm,bottom=2.6mm,before skip=11pt,
  after={\par\nopagebreak\vskip7pt}]
{\sffamily\bfseries\fontsize{15}{18}\selectfont\milctintasobre{#1}#2}
\end{tcolorbox}}

% Tarjeta de paso numerada: sustituye las tablas «Pilar / Aplicacion», que
% eran formato de consulta docente, no de estudiante. El numero va en un
% circulo sobre el borde para que la secuencia se lea de un vistazo.
\newcommand{\milcpaso}[3]{%
\Needspace{4\baselineskip}%
\begin{tcolorbox}[enhanced,colback=white,colframe=milcLinea,arc=1.5mm,boxrule=.9pt,
  left=14mm,right=4mm,top=3mm,bottom=3mm,before skip=4pt,after skip=4pt,
  fontupper=\RaggedRight,
  overlay={\node[anchor=north west,circle,fill=#2,text=white,inner sep=0pt,
    minimum size=7.4mm,font=\sffamily\bfseries\small]
    at ([xshift=3.6mm,yshift=-3.2mm]frame.north west) {#1};}]
#3
\end{tcolorbox}}

% Casilla marcable de tamano util con lapiz (la anterior era \fbox{\phantom{x}},
% de alto variable segun el texto que la seguia).
\newcommand{\milcmarca}{\framebox[3.8mm]{\rule{0pt}{3.4mm}}}

% Fila marcable de lista suelta (checks de la estacion 1).
\newcommand{\milccheckrow}[1]{%
\par\vspace{1.8mm}\noindent
\makebox[6.4mm][l]{\raisebox{-.55mm}{\textcolor{milcNegro}{\milcmarca}}}%
\parbox[t]{\dimexpr\linewidth-6.4mm\relax}{\RaggedRight #1}\par}

% Celda marcable dentro de la rubrica (tabla de autoevaluacion). Misma
% sangria francesa, pero pensada para vivir dentro de una columna X.
\newcommand{\milcopcion}[1]{%
\makebox[5.2mm][l]{\raisebox{-.55mm}{\textcolor{milcNegro}{\milcmarca}}}%
\parbox[t]{\dimexpr\linewidth-5.2mm\relax}{\RaggedRight #1}}

% ── Recursos visuales de la guía ──────────────────────────────────────
% Declarados en el bloque `recursos:` del YAML y emitidos por el builder.
% \guiaFigura   → imagen rasterizada o PDF (foto, carta celeste, montaje).
% \guiaDiagrama → fuente TikZ/circuitikz (.tex) incrustada como vector:
%                 es la vía para circuitos y esquemáticos editables.
% [#1] = texto alternativo (alt) · #2 = ruta relativa al .tex · #3 = pie
% El alt llega al PDF etiquetado (/Alt) y lo lee el lector de pantalla.
\newcommand{\guiaFigura}[3][]{%
\begin{center}
\includegraphics[width=0.88\linewidth,height=0.38\textheight,keepaspectratio,alt={#1}]{#2}\\[1.5mm]
{\footnotesize\itshape\textcolor{milcNegro!75}{#3}}
\end{center}
}

\newcommand{\guiaDiagrama}[2]{%
\begin{center}
\input{#1}\\[1.5mm]
{\footnotesize\itshape\textcolor{milcNegro!75}{#2}}
\end{center}
}

\begin{document}
\thispagestyle{empty}

% ── Portada ───────────────────────────────────────────────────────────
% Antes: un rectangulo de color a pagina completa. Se veia bien en pantalla,
% pero en la fotocopiadora del colegio se comia el toner y salia gris sucio.
% Ahora la banda ocupa el tercio superior y el resto va en blanco.
\begin{tikzpicture}[remember picture,overlay]
  \path[fill=milcGradePrimary]
    (current page.north west) rectangle ([yshift=-10.6cm]current page.north east);

  \node[anchor=north west,text=milcGradeText] at ([xshift=2.0cm,yshift=-1.55cm]current page.north west)
    {\sffamily\bfseries\fontsize{12}{14}\selectfont GRADO};
  \node[anchor=north west,text=milcGradeText,opacity=.85] at ([xshift=2.0cm,yshift=-2.35cm]current page.north west)
    {\sffamily\bfseries\fontsize{8.5}{10}\selectfont SERIE GUÍAS · TECNOLOGÍA E INFORMÁTICA};
  \node[anchor=north east,text=milcGradeText,opacity=.85,align=right] at ([xshift=-2.0cm,yshift=-1.55cm]current page.north east)
    {\sffamily\bfseries\fontsize{9}{12}\selectfont I.E. Sor María Juliana\\Cartago, Valle del Cauca};

  \node[anchor=west,text=milcGradeText] (gradonum) at ([xshift=1.85cm,yshift=-7.1cm]current page.north west)
    {\sffamily\bfseries\fontsize{112}{112}\selectfont 9};
  % El circulito de grado (°) solo aplica a las guias de grado («11°»); en
  % Bebras y Semillero el numero grande es un momento/modulo, no un grado.
  % Se posiciona relativo al nodo `gradonum`, no en coordenadas absolutas.
  \draw[milcGradeText,line width=1.6pt]
    ([xshift=2.0mm,yshift=-4.5mm]gradonum.north east) circle (.15cm);

  \node[anchor=west,text=milcGradeText,text width=11.4cm,align=left] at ([xshift=1.5cm,yshift=0.5cm]gradonum.east)
    {
     {\sffamily\bfseries\fontsize{9}{11}\selectfont\MakeUppercase{Período 1 · Historia de la técnica}}\\[3mm]
     {\sffamily\bfseries\fontsize{21}{25}\selectfont\setlength{\baselineskip}{27pt}La rueda y el agua ---\\[4pt]máquinas simples\\[4pt]del campo}\\[3.5mm]
     {\sffamily\bfseries\fontsize{15}{18}\selectfont Guía 2-9-TIC}
    };
\end{tikzpicture}

\vspace*{9.4cm}
\vfill

{\sffamily\bfseries\fontsize{8}{10}\selectfont\textcolor{milcGradeDark}{LO QUE TE LLEVAS HOY}}

\begin{tcolorbox}[enhanced,colback=white,colframe=milcNegro,arc=2mm,boxrule=1.6pt,
  left=5mm,right=5mm,top=4mm,bottom=4mm,before skip=3pt,after skip=12pt,
  fontupper=\RaggedRight\fontsize{11}{15.5}\selectfont]
Tres dibujos a mano de máquinas del campo colombiano: un trapiche, una noria o molino de agua y una polea. Cada dibujo lleva al menos cuatro piezas rotuladas con su nombre, dice cuál de las seis máquinas simples opera adentro y lo justifica en una frase. Cada uno responde además las dos preguntas que importan: qué multiplica y qué pide a cambio.
\end{tcolorbox}

\begin{tcolorbox}[enhanced,colback=milcGris,colframe=milcGris,arc=2mm,boxrule=0pt,
  left=5mm,right=5mm,top=3.5mm,bottom=3.5mm,after skip=12pt]
{\sffamily\bfseries\fontsize{8}{10}\selectfont\textcolor{milcNegro!55}{ESTA GUÍA ES DE}}\\[3mm]
\begin{tabularx}{\linewidth}{X p{3.2cm} p{3.2cm}}
\linead{\linewidth} & \linead{3.0cm} & \linead{3.0cm}\\[-1mm]
{\fontsize{8.5}{10}\selectfont\textcolor{milcNegro!55}{Nombre completo}} &
{\fontsize{8.5}{10}\selectfont\textcolor{milcNegro!55}{Grupo}} &
{\fontsize{8.5}{10}\selectfont\textcolor{milcNegro!55}{Fecha}}\\
\end{tabularx}
\end{tcolorbox}

\vfill

\noindent\textcolor{milcLinea}{\rule{\linewidth}{1pt}}\\[2mm]
{\sffamily\bfseries\fontsize{9.5}{12}\selectfont Autor: Dr. Álvaro Cárdenas Orozco}\\[1mm]
{\sffamily\fontsize{8.5}{11}\selectfont\textcolor{milcNegro!55}{Metodología MILC · apertura ancestral · pensamiento computacional · triángulo Dussel–estoicismo–Floridi}}

\newpage

\section*{Apertura: La rueda y el agua --- máquinas simples del campo colombiano}

\begin{softbox}{milcGradeDark}{milcGradeSoft}{Saber ancestral}
El río La Vieja bordea a Cartago y le pone frontera con Risaralda. Ese mismo río movió trapiches y molinos, porque el agua que corre es fuerza gratis para quien sepa tomarla. Pero el agua cobra. Cuando crece no se lleva la ciudad entera: se lleva siempre los mismos barrios. Brisas del Río, La Platanera, La Arenera, La Playa y Tierra del Olvido quedaron bajo el agua en marzo de 2022. El río llegó a 12,3 metros. Y la gente no adivina la creciente: la mira. \textbf{Puerto Alejandría} es el punto al que se va a ver el nivel. Es la referencia con la que se compara cómo está hoy contra cómo estaba ayer. Ese es el miedo que protege: no el que paraliza, sino el que enseña a qué hay que mirarle el nivel. La \textbf{cara de exclusión}: el riesgo está repartido por clase y por ubicación, no por azar.\par\smallskip{\footnotesize\textcolor{milcNegro!70}{\textbf{Fuente.} CiudadRegión. (2022, 7 de marzo). \emph{Alarma por alto nivel del río La Vieja a su paso por Cartago}. CiudadRegión Noticias.}}
\end{softbox}

\begin{softbox}{milcTurquesa}{milcGris}{Saber contemporáneo}
Hay seis máquinas simples y todas las demás son combinaciones de ellas. \textbf{Palanca}, \textbf{rueda con eje}, \textbf{polea}, \textbf{plano inclinado}, \textbf{cuña} y \textbf{tornillo}. Cada una hace lo mismo en el fondo: cambia la relación entre la fuerza que pones y el efecto que consigues. Eso se llama \textbf{ventaja mecánica}. Y aquí está la regla que nadie puede saltarse: ninguna máquina regala energía. Si ganas fuerza, pierdes distancia. Si ganas distancia o velocidad, pierdes fuerza. El trapiche multiplica lo que empuja el buey, y a cambio el buey camina muchas vueltas. La polea del pozo te deja subir el balde tirando hacia abajo, y a cambio jalas la misma cantidad de cuerda. Por eso una sola pregunta basta para entender cualquier máquina: qué multiplica y qué pide a cambio. Si solo sabes contestar la primera mitad, todavía no la entendiste.
\end{softbox}

\begin{softbox}{milcMostaza}{milcGris}{Pregunta-puente}
\textit{El río da fuerza gratis y cobra creciendo sobre los mismos barrios. ¿Qué máquina de tu casa te da algo a cambio de un costo que casi nunca miras?}
\end{softbox}

\begin{softbox}{milcVerde}{milcGris}{Saber y saber hacer}
Al terminar podrás: (1) \textbf{identificar} las seis máquinas simples en objetos de tu entorno, con un ejemplo de cada una; (2) \textbf{explicar} qué multiplica cada máquina y qué pide a cambio, sin saltarte la segunda mitad; (3) \textbf{crear} tres dibujos analíticos de máquinas del campo colombiano con sus piezas rotuladas.
\end{softbox}



\small
\begin{tabularx}{\textwidth}{>{\bfseries}p{3.0cm}Y}
\rowcolor{milcGradePrimary!12}Autor & Dr. Álvaro Cárdenas Orozco\\
DBA & Análisis crítico de la historia de la tecnología y su impacto social (MEN, T\&I)\\
Referentes & Dussel (técnica situada) · Estoicismo (téchne del cuerpo) · Floridi (infoesfera)\\
Producto final & Tres dibujos a mano de máquinas del campo colombiano: un trapiche, una noria o molino de agua y una polea. Cada dibujo lleva al menos cuatro piezas rotuladas con su nombre, dice cuál de las seis máquinas simples opera adentro y lo justifica en una frase. Cada uno responde además las dos preguntas que importan: qué multiplica y qué pide a cambio.\\
\end{tabularx}
\normalsize

\puente{En la sesión 1 definiste qué es técnica y viste que la máquina automatiza con energía propia. Hoy abres esa caja y miras qué hay adentro. En la sesión 3 llega la balanza, que es una palanca dedicada a medir.}

\milcestacion{milcGradeDark}{Ruta MILC: así vamos a aprender}
\begin{tabularx}{\textwidth}{>{\centering\arraybackslash}X>{\centering\arraybackslash}X>{\centering\arraybackslash}X>{\centering\arraybackslash}X}
\phasecard{1}{milcVerde}{milcGris}{Escucha\\[-1mm]\small Observo, escucho y pregunto.} &
\phasecard{2}{milcTurquesa}{milcGris}{Entender\\[-1mm]\small Sistematizo y ordeno.} &
\phasecard{3}{milcMagenta}{milcGris}{Hacer\\[-1mm]\small Construyo y aplico.} &
\phasecard{4}{milcVino}{milcGris}{Cierre\\[-1mm]\small Evalúo y reflexiono.}
\end{tabularx}


\puente{Antes de la teoría vas a cazar las seis máquinas simples a tu alrededor. Están todas, aunque no las llames así.}

\milcestacion{milcVerde}{Estación 1 de 3 · Escucha}

\begin{softbox}{milcVerde}{milcGris}{Escena para escuchar}
\textbf{Actividad 1 · IDENTIFICA --- Las seis a tu alrededor} (15 min · individual). (1) Encuentra un ejemplo cotidiano de cada una de las seis máquinas simples. Palanca, rueda con eje, polea, plano inclinado, cuña y tornillo. (2) Pueden estar en tu casa, tu mochila, la calle o una bicicleta. (3) Escribe para cada una qué hace y qué te ahorra. (4) Marca las dos que te costó más encontrar y anota por qué.
\end{softbox}

{\sffamily\bfseries\fontsize{8}{10}\selectfont\textcolor{milcNegro!55}{MARCA CUANDO LO HAYAS HECHO}}
\milccheckrow{Tengo un ejemplo real de cada una de las seis.}
\milccheckrow{Escribí qué hace y qué me ahorra en cada caso.}
\milccheckrow{Marqué las dos que me costó más encontrar.}

\begin{infoband}{milcVerde}


\textbf{Lo que va al cuaderno (Actividad 1 · IDENTIFICA).} \textbf{Título:} «Actividad 1 --- Las seis a mi alrededor». \textbf{Formato:} tabla de 6 filas y 3 columnas (máquina simple / objeto real / qué me ahorra), con marca en las dos más difíciles. \textbf{Extensión:} un tercio de página. \textbf{Sabes que terminaste cuando} las seis filas tienen un objeto concreto, no una categoría.
\end{infoband}

\puente{Ya las encontraste. Ahora vas a entender por qué ninguna regala nada, y a poder decir en cada caso qué se paga.}

\milcestacion{milcTurquesa}{Estación 2 de 3 · Entender}

\begin{softbox}{milcTurquesa}{milcGris}{Lo que vas a entender}
\textbf{Actividad 2 · EXPLICA --- Qué multiplica y qué pide} (20 min · parejas). (1) Con tu pareja, tomen tres de sus doce objetos y escriban quién aplica la fuerza en cada uno. (2) Escriban qué multiplica cada uno: fuerza, velocidad o dirección. (3) Escriban qué pide a cambio, que siempre es distancia, tiempo o recorrido. (4) Busquen un objeto donde la segunda parte no sea evidente y discútanlo hasta poder escribirla.
\end{softbox}

\milcpaso{1}{milcTurquesa}{\textbf{Las seis y ninguna más.} Palanca, rueda con eje, polea, plano inclinado, cuña y tornillo. Todo lo demás es combinación. Un trapiche es rueda y palanca; una escalera es plano inclinado repetido.}
\milcpaso{2}{milcTurquesa}{\textbf{Ninguna regala energía.} Si ganas fuerza, pierdes distancia. Si ganas distancia, pierdes fuerza. Esa es la regla, y no depende del material ni de la época. El buey del trapiche camina muchas vueltas por una razón física, no por costumbre.}
\milcpaso{3}{milcTurquesa}{\textbf{La pregunta doble.} Qué multiplica y qué pide a cambio. Casi todo el mundo sabe contestar la primera mitad. La segunda es la que separa entender de repetir.}
\milcpaso{4}{milcTurquesa}{\textbf{«Simple» no es «primitivo».} Simple quiere decir que no se puede descomponer en algo más básico. Una noria bien hecha aprovecha una fuerza gratuita durante décadas sin consumir nada, y eso no lo hace ningún motor.}

\begin{drawbox}{milcTurquesa}{Cualquier máquina, en cuatro preguntas}
\textbf{1. Fuerza:} ¿quién la aplica, persona, animal o agua? \\[1mm] \textbf{2. Máquina simple:} ¿cuál de las seis opera adentro? \\[1mm] \textbf{3. Multiplica:} ¿fuerza, velocidad o dirección? \\[1mm] \textbf{4. Pide a cambio:} ¿distancia, tiempo o recorrido?
\end{drawbox}

\begin{softbox}{milcMostaza}{milcGris}{Errores comunes y su corrección}
(1) \textbf{Creer que la máquina crea fuerza}: la redistribuye, no la fabrica. (2) \textbf{Contestar solo qué multiplica}: sin el costo, el análisis está a medias. (3) \textbf{Confundir simple con primitivo}: simple es que no se descompone más. (4) \textbf{Olvidar quién aplica la fuerza}: en el trapiche es el buey o el motor, y eso cambia todo el análisis. (5) \textbf{Dibujar sin rotular}: un dibujo sin nombres no se puede leer ni corregir.
\end{softbox}

\begin{infoband}{milcTurquesa}


\textbf{Lo que va al cuaderno (Actividad 2 · EXPLICA).} \textbf{Título:} «Actividad 2 --- Qué multiplica y qué pide». \textbf{Formato:} tres objetos analizados con las cuatro preguntas, más el caso difícil con su costo escrito. \textbf{Extensión:} media página. \textbf{Sabes que terminaste cuando} los tres tienen escrita la segunda mitad, la del costo.
\end{infoband}

\puente{Tienes la regla. Toca dibujar tres máquinas reales del campo y señalar dentro de cada una cuál de las seis está operando.}

\milcestacion{milcMagenta}{Estación 3 de 3 · Hacer}

\begin{softbox}{milcMagenta}{milcGris}{Cómo vas a construir tu producto}
\textbf{Actividad 3 · CREA --- Tres dibujos analíticos} (30 min · individual). (1) Dibuja a mano un trapiche, una noria o molino de agua y una polea. (2) Rotula al menos cuatro piezas de cada uno con su nombre. (3) Escribe en cada dibujo cuál de las seis máquinas simples opera adentro. (4) Justifica esa identificación en una frase. (5) Añade en cada uno qué multiplica y qué pide a cambio. (6) Pásaselos a un compañero y anota qué no entendió sin que le explicaras.
\end{softbox}

\milcpaso{1}{milcMagenta}{\textbf{Tu entrega tiene tres dibujos y una prueba.} Los tres con piezas rotuladas, máquina simple identificada y la pregunta doble contestada. Y lo que dijo tu compañero al leerlos solo.}
\milcpaso{2}{milcMagenta}{\textbf{Dibujar a mano obliga a entender.} Copiar una foto no sirve: al dibujar hay que decidir qué pieza va sobre cuál, y ahí se ve si sabes cómo funciona. Los dibujos feos que están bien valen más que los bonitos que están mal.}
\milcpaso{3}{milcMagenta}{\textbf{Rotular es la mitad del trabajo.} «Eje», «canal», «cangilón», «masa», «cuja». Si no conoces el nombre de una pieza, averígualo o descríbela; lo que no vale es dejarla sin nombre.}
\milcpaso{4}{milcMagenta}{\textbf{La noria es el caso más interesante.} Usa una fuerza que no se agota y que no le cuesta a nadie... salvo cuando el río crece. Ese doble filo es exactamente lo que hay que anotar.}

\begin{drawbox}{milcMagenta}{Antes de entregar}
\checkbox Tres dibujos a mano: trapiche, noria o molino, y polea. \\[1mm] \checkbox Al menos cuatro piezas rotuladas en cada uno. \\[1mm] \checkbox Identificada la máquina simple que opera adentro. \\[1mm] \checkbox Justificación en una frase por dibujo. \\[1mm] \checkbox Qué multiplica y qué pide a cambio, en los tres. \\[1mm] \checkbox Un compañero los leyó sin explicación.
\end{drawbox}

\begin{softbox}{milcMagenta}{milcGris}{Plantilla de trabajo}
\textbf{Máquina.} \textit{Cuál es y dónde se usa.} \\[1mm] \textbf{Piezas.} \textit{Cuatro nombres rotulados.} \\[1mm] \textbf{Máquina simple.} \textit{Cuál de las seis y por qué.} \\[1mm] \textbf{Multiplica.} \textit{Fuerza, velocidad o dirección.} \\[1mm] \textbf{Pide a cambio.} \textit{Distancia, tiempo o recorrido.}
\end{softbox}

\begin{infoband}{milcMagenta}


\textbf{Lo que va al cuaderno (Actividad 3 · CREA).} \textbf{Título:} «Actividad 3 --- Tres dibujos analíticos». \textbf{Formato:} los tres dibujos con sus piezas rotuladas, la máquina simple identificada y la pregunta doble contestada en cada uno. \textbf{Extensión:} una página. \textbf{Sabes que terminaste cuando} un compañero leyó los tres sin preguntarte nada.
\end{infoband}

\puente{Lo que entregas hoy: tres dibujos con piezas rotuladas y la máquina simple identificada. La prueba de calidad: un compañero los lee sin que le expliques nada.}

\milcestacion{milcMostaza}{Producto de la sesión}
\textbf{Tres dibujos analíticos de máquinas del campo con piezas rotuladas}.
\begin{softbox}{milcMostaza}{milcGris}{Criterios de calidad}
(1) Los tres dibujos están hechos a mano, no copiados de una foto. (2) Cada uno tiene al menos cuatro piezas rotuladas con su nombre. (3) Cada uno identifica cuál de las seis máquinas simples opera adentro. (4) Cada uno dice qué multiplica y qué pide a cambio. (5) Un compañero los leyó sin explicación adicional.
\end{softbox}

\puente{Antes de cerrar, mira lo que hiciste desde cinco lados. «Simple» no quiere decir «primitivo»: quiere decir que no se puede descomponer en algo más básico.}

\milcestacion{milcVino}{Evaluación liberadora · cinco dimensiones}

\begin{softbox}{milcVino}{milcGris}{Sentido de la evaluación}
Evaluar no es solo revisar una nota. Es reconocer qué aprendí, qué cambió en mí, qué emociones manejé, cómo me relaciono con la ciudad y la comunidad, y qué le dejaré a quien viene detrás.
\end{softbox}

\textbf{Mi reflexión liberadora en cinco dimensiones.} Completa cada frase iniciada:

\small
\begin{tabularx}{\textwidth}{>{\bfseries\RaggedRight\arraybackslash}p{3.4cm}Y>{\RaggedRight\arraybackslash}p{4.6cm}}
\rowcolor{milcVino}\color{white}Dimensión & \color{white}\bfseries Mi respuesta (frame starter) & \color{white}\bfseries Algo que descubrí\\
1. Personal & Lo que más cambió en mí fue \linead{6.5cm} & \linead{4.2cm}\\
2. Emocional & La emoción más fuerte fue \linead{2.5cm} y la regulé \linead{3cm} & \linead{4.2cm}\\
3. Ciudadana & En democracia esto importa porque \linead{6cm} & \linead{4.2cm}\\
4. Local (Cartago) & En mi barrio sería útil para \linead{6.5cm} & \linead{4.2cm}\\
5. Intergeneracional & A mi abuela/abuelo le diría \linead{6.5cm} & \linead{4.2cm}\\
\end{tabularx}
\normalsize

\begin{infoband}{milcVino}
\textbf{Para la próxima sustentación.} Vuelve a esta tabla en seis meses. Verás que las respuestas cambian — no porque lo aprendido cambie, sino porque tú cambias. La evaluación liberadora es una conversación contigo a través del tiempo.
\end{infoband}


\puente{Para terminar, tres citas verificadas sobre diseño y costo: Dussel sobre la coherencia formal del producto, Séneca sobre lo que nos aterra frente a lo que nos aprieta, Floridi sobre lo único que no se puede compartir.}

\milcestacion{milcGradeDark}{Triángulo de pensamiento · cierre filosófico}

\begin{softbox}{milcVino}{milcGris}{Dussel · lente del nosotros}
\textit{«El tema esencial del diseño es el de dotar al producto de coherencia formal; incluye a la tecnología… por cuanto esto significa coherencia funcional: el del valor del uso; incluye a la estética, porque la coherencia formal, en cuanto tal, es la belleza del producto.»}

{\footnotesize\itshape\color{milcNegro!70}--- Enrique Dussel · Filosofía de la liberación (1977), §4.3.1.2}

Un trapiche bien hecho es bonito porque está bien resuelto, no al revés. Cuando dibujes, fíjate en que casi ninguna pieza sobra: eso es coherencia formal, y se nota en objetos que nadie diseñó en una oficina.

\textbf{¿Qué objeto de mi casa me parece bonito solo porque está bien resuelto?}
\end{softbox}
\linead{\linewidth}\\[1mm]
\linead{\linewidth}

\begin{softbox}{milcMostaza}{milcGris}{Estoicismo · Séneca · Cartas a Lucilio, 13 (c. 64 d.C.) · cuidado interior}
\textit{«Muchas más son, Lucilio, las cosas que nos aterran que las que realmente nos aprietan; frecuentemente sufrimos más por las opiniones que por la realidad.»}

{\footnotesize\itshape\color{milcNegro!70}--- Séneca · Cartas a Lucilio, 13 (c. 64 d.C.)}

En Cartago la gente no adivina la creciente: mira el nivel en Puerto Alejandría. Esa es la diferencia entre el miedo que paraliza y el que protege. Un dato medido tranquiliza más que cien suposiciones.

\textbf{¿Qué cosa me preocupa sin que yo haya mirado nunca su nivel de verdad?}
\end{softbox}
\linead{\linewidth}\\[1mm]
\linead{\linewidth}

\begin{softbox}{milcTurquesa}{milcGris}{Floridi · lente de la infoesfera}
\textit{«Lo que en última instancia es finito, precioso, no renovable e incompartible es, en realidad, el tiempo. (trad. propia)»}

{\footnotesize\itshape\color{milcNegro!70}--- Luciano Floridi · Commentary on the Onlife Manifesto (2015), § 4.6}

Esa es la moneda con la que se paga casi toda ventaja mecánica. La polea te ahorra fuerza y te cobra cuerda; el trapiche ahorra esfuerzo y cobra vueltas. Casi siempre lo que se entrega a cambio es tiempo.

\textbf{¿Qué máquina uso a diario que me ahorra esfuerzo y me cobra tiempo sin que lo note?}
\end{softbox}
\linead{\linewidth}\\[1mm]
\linead{\linewidth}

\begin{infoband}{milcNegro}
\textbf{Nota para el docente.} Las tres son citas textuales y llevan su referencia debajo. Puedes remitir a la obra en clase.
\end{infoband}

\milcestacion{milcVino}{Mi autoevaluación · marca una casilla por fila}

{\small
\setlength{\extrarowheight}{2.2pt}
\begin{tabularx}{\textwidth}{>{\RaggedRight\arraybackslash\bfseries}p{3.0cm}YYY}
\rowcolor{milcVino}\color{white}\bfseries Criterio & \color{white}\bfseries Lo logré & \color{white}\bfseries En proceso & \color{white}\bfseries Necesito apoyo\\[.6mm]
Comprensión del tema & \milcopcion{Explico las ideas con mis palabras.} & \milcopcion{Reconozco las ideas, falta precisión.} & \milcopcion{Necesito repasar lo básico.}\\[.8mm]
\rowcolor{milcGris}Pensamiento computacional & \milcopcion{Usé los pilares al armar mi producto.} & \milcopcion{Usé algunos pilares.} & \milcopcion{Necesito releer la estación 2.}\\[.8mm]
Aplicación y producto & \milcopcion{Mi producto cumple los criterios.} & \milcopcion{Está armado, con ajustes pendientes.} & \milcopcion{Aún está en construcción.}\\[.8mm]
\rowcolor{milcGris}Reflexión 5 dimensiones & \milcopcion{Respondí cada una con honestidad.} & \milcopcion{Respondí algunas con detalle.} & \milcopcion{Mis respuestas son muy generales.}\\[.8mm]
Triángulo de pensamiento & \milcopcion{Conecté las 3 voces con mi trabajo.} & \milcopcion{Conecté al menos una.} & \milcopcion{No las relaciono con mi proceso.}\\[.8mm]
\end{tabularx}}

\vspace{3mm}
\textbf{Hoy entendí que:}\\[1mm]
\linead{\linewidth}\\[1mm]
\linead{\linewidth}

\textbf{Mi compromiso para la próxima sesión es:}\\[1mm]
\linead{\linewidth}\\[1mm]
\linead{\linewidth}

\begin{softbox}{milcMostaza}{milcGris}{Cita de despedida · Marco Aurelio}
\textit{``El obstáculo en el camino se vuelve el camino. Lo que estorba al camino se vuelve camino.''}\\[1mm]
Cada error de hoy, bien observado, se convierte en pieza del camino que pisarás mañana.
\end{softbox}

\small
\begin{tabularx}{\textwidth}{>{\bfseries}p{3.6cm}Y}
\rowcolor{milcGradeDark!12}Nombre completo: & \linead{10cm}\\
Fecha: & \linead{4cm} \quad Curso/grupo: \linead{4cm}\\
Mi firma: & \linead{9cm}\\
\end{tabularx}
\normalsize

\begin{infoband}{milcGradeDark}
\textbf{Para la próxima sesión.} Llega con los tres dibujos rotulados. Si puedes, mira una máquina real de tu barrio y cuéntale a alguien qué multiplica y qué pide a cambio. En la sesión 3 vas a construir una balanza, que es una palanca dedicada a medir.
\end{infoband}

\end{document}
